% =====================================================================
%  reporte.tex
%  ========================
%
%
%  Descripción:
%  ------------
%  Reporte de la actividad 1 de Fundamentos para el Análisis Automático
%  de Lenguaje: preprocesamiento de la colección Time y extracción del
%  vocabulario con su frecuencia de término.
%
%
%  Metadata:
%  ----------
%   * Author : zxxz6 (Bryan Violante Arriaga)
%   * Version: 1.5.0
%
%
%  History:
%  ------------
%  Author      Date            Description
%  zxxz6       09/09/2026      Creation
%
%
% =====================================================================
\documentclass[11pt,letterpaper]{article}

% =====================================================================
%  DATOS DE LA MATERIA
% =====================================================================
\newcommand{\materia}{Fundamentos para el Análisis Automático de Lenguaje}
\newcommand{\profesor}{Dra. Delia Irazú Hernández Farías \\
                       Dr. Luis Villaseñor Pineda \\
                       Dr. Manuel Montes y Gómez \\
                       Dr. Aurelio López López}
\newcommand{\periodo}{Otoño 2026}
\newcommand{\alumno}{Bryan Ezequiel Violante Arriaga}

% =====================================================================
%  Preambulo.tex
%  ========================
%
%
%  Descripción:
%  ------------
%  Preámbulo compartido de los apuntes de clase: paquetes, estilo de
%  página, entornos de teorema, las cuatro cajas de color, el estilo de
%  listings y el pseudocódigo en español. Se carga con \input desde
%  cada documento en lugar de copiarlo, para que un arreglo aquí llegue
%  a todas las libretas de una vez.
%
%
%  Considerations:
%  ------------
%  - Se carga con \input y ruta relativa, no con \usepackage, porque
%    \usepackage solo busca en la carpeta del documento y en el árbol
%    de TeX. Un .sty instalado en TEXMFHOME evitaría la ruta, pero vive
%    fuera del repo y una clona nueva no compilaría
%  - Los cuatro datos de la materia se declaran en el documento ANTES
%    del \input. hyperref expande \materia al cargarse, así que si no
%    existe todavía la compilación truena
%  - Cuando la materia la imparte más de un profesor, los nombres van
%    en un solo \profesor separados por \\. La portadilla los imprime
%    dentro de un center, que es donde \\ sí corta renglón. Repetir
%    \newcommand{\profesor} no es opción: truena con "already defined"
%  - algorithmicx no tiene opción de idioma y babel no lo traduce, así
%    que las palabras clave del pseudocódigo se renombran a mano
%  - Los nombres de los entornos van sin acento (\begin{definicion})
%    aunque el rótulo impreso sí lo lleve
%  - babel se carga con es-noshorthands porque en español el " queda
%    activo y se come el texto que le sigue, además de romper listings
%
%
%  Metadata:
%  ----------
%   * Author : zxxz6 (Bryan Violante Arriaga)
%   * Version: 1.7.0
%
%
%  History:
%  ------------
%  Author      Date            Description
%  zxxz6       24/09/2026      El nombre del alumno en el encabezado
%  zxxz6       24/09/2026      Encabezado de una hoja de un ejercicio
%  zxxz6       23/09/2026      Macros de los documentos de soluciones
%  zxxz6       30/08/2026      Macros de lim, limsup y liminf sobre n
%  zxxz6       25/08/2026      Macros de Omega, Theta, o y omega chicas
%  zxxz6       20/08/2026      Cargue pgfplots para graficas de datos
%  zxxz6       20/08/2026      Cargue tikz para diagramas de conjuntos
%  zxxz6       19/08/2026      Documente \profesor con varios nombres
%  zxxz6       19/08/2026      Agregue la caja Idea Random
%  zxxz6       18/08/2026      Creation
%
%
% =====================================================================
%  USO
%  ---
%  \documentclass[11pt,letterpaper]{article}
%  \newcommand{\materia}{...}      % los cuatro datos van arriba
%  \newcommand{\profesor}{...}
%  \newcommand{\periodo}{...}
%  \newcommand{\alumno}{...}
%  % =====================================================================
%  Preambulo.tex
%  ========================
%
%
%  Descripción:
%  ------------
%  Preámbulo compartido de los apuntes de clase: paquetes, estilo de
%  página, entornos de teorema, las cuatro cajas de color, el estilo de
%  listings y el pseudocódigo en español. Se carga con \input desde
%  cada documento en lugar de copiarlo, para que un arreglo aquí llegue
%  a todas las libretas de una vez.
%
%
%  Considerations:
%  ------------
%  - Se carga con \input y ruta relativa, no con \usepackage, porque
%    \usepackage solo busca en la carpeta del documento y en el árbol
%    de TeX. Un .sty instalado en TEXMFHOME evitaría la ruta, pero vive
%    fuera del repo y una clona nueva no compilaría
%  - Los cuatro datos de la materia se declaran en el documento ANTES
%    del \input. hyperref expande \materia al cargarse, así que si no
%    existe todavía la compilación truena
%  - Cuando la materia la imparte más de un profesor, los nombres van
%    en un solo \profesor separados por \\. La portadilla los imprime
%    dentro de un center, que es donde \\ sí corta renglón. Repetir
%    \newcommand{\profesor} no es opción: truena con "already defined"
%  - algorithmicx no tiene opción de idioma y babel no lo traduce, así
%    que las palabras clave del pseudocódigo se renombran a mano
%  - Los nombres de los entornos van sin acento (\begin{definicion})
%    aunque el rótulo impreso sí lo lleve
%  - babel se carga con es-noshorthands porque en español el " queda
%    activo y se come el texto que le sigue, además de romper listings
%
%
%  Metadata:
%  ----------
%   * Author : zxxz6 (Bryan Violante Arriaga)
%   * Version: 1.7.0
%
%
%  History:
%  ------------
%  Author      Date            Description
%  zxxz6       24/09/2026      El nombre del alumno en el encabezado
%  zxxz6       24/09/2026      Encabezado de una hoja de un ejercicio
%  zxxz6       23/09/2026      Macros de los documentos de soluciones
%  zxxz6       30/08/2026      Macros de lim, limsup y liminf sobre n
%  zxxz6       25/08/2026      Macros de Omega, Theta, o y omega chicas
%  zxxz6       20/08/2026      Cargue pgfplots para graficas de datos
%  zxxz6       20/08/2026      Cargue tikz para diagramas de conjuntos
%  zxxz6       19/08/2026      Documente \profesor con varios nombres
%  zxxz6       19/08/2026      Agregue la caja Idea Random
%  zxxz6       18/08/2026      Creation
%
%
% =====================================================================
%  USO
%  ---
%  \documentclass[11pt,letterpaper]{article}
%  \newcommand{\materia}{...}      % los cuatro datos van arriba
%  \newcommand{\profesor}{...}
%  \newcommand{\periodo}{...}
%  \newcommand{\alumno}{...}
%  % =====================================================================
%  Preambulo.tex
%  ========================
%
%
%  Descripción:
%  ------------
%  Preámbulo compartido de los apuntes de clase: paquetes, estilo de
%  página, entornos de teorema, las cuatro cajas de color, el estilo de
%  listings y el pseudocódigo en español. Se carga con \input desde
%  cada documento en lugar de copiarlo, para que un arreglo aquí llegue
%  a todas las libretas de una vez.
%
%
%  Considerations:
%  ------------
%  - Se carga con \input y ruta relativa, no con \usepackage, porque
%    \usepackage solo busca en la carpeta del documento y en el árbol
%    de TeX. Un .sty instalado en TEXMFHOME evitaría la ruta, pero vive
%    fuera del repo y una clona nueva no compilaría
%  - Los cuatro datos de la materia se declaran en el documento ANTES
%    del \input. hyperref expande \materia al cargarse, así que si no
%    existe todavía la compilación truena
%  - Cuando la materia la imparte más de un profesor, los nombres van
%    en un solo \profesor separados por \\. La portadilla los imprime
%    dentro de un center, que es donde \\ sí corta renglón. Repetir
%    \newcommand{\profesor} no es opción: truena con "already defined"
%  - algorithmicx no tiene opción de idioma y babel no lo traduce, así
%    que las palabras clave del pseudocódigo se renombran a mano
%  - Los nombres de los entornos van sin acento (\begin{definicion})
%    aunque el rótulo impreso sí lo lleve
%  - babel se carga con es-noshorthands porque en español el " queda
%    activo y se come el texto que le sigue, además de romper listings
%
%
%  Metadata:
%  ----------
%   * Author : zxxz6 (Bryan Violante Arriaga)
%   * Version: 1.7.0
%
%
%  History:
%  ------------
%  Author      Date            Description
%  zxxz6       24/09/2026      El nombre del alumno en el encabezado
%  zxxz6       24/09/2026      Encabezado de una hoja de un ejercicio
%  zxxz6       23/09/2026      Macros de los documentos de soluciones
%  zxxz6       30/08/2026      Macros de lim, limsup y liminf sobre n
%  zxxz6       25/08/2026      Macros de Omega, Theta, o y omega chicas
%  zxxz6       20/08/2026      Cargue pgfplots para graficas de datos
%  zxxz6       20/08/2026      Cargue tikz para diagramas de conjuntos
%  zxxz6       19/08/2026      Documente \profesor con varios nombres
%  zxxz6       19/08/2026      Agregue la caja Idea Random
%  zxxz6       18/08/2026      Creation
%
%
% =====================================================================
%  USO
%  ---
%  \documentclass[11pt,letterpaper]{article}
%  \newcommand{\materia}{...}      % los cuatro datos van arriba
%  \newcommand{\profesor}{...}
%  \newcommand{\periodo}{...}
%  \newcommand{\alumno}{...}
%  \input{../../templates/Preambulo.tex}
%  \begin{document}
%  \portadilla
%
%  Con dos o mas profesores se declara UN solo \profesor y los nombres
%  se separan con \\, uno por renglon:
%
%  \newcommand{\profesor}{Dra. Fulana de Tal \\
%                         Dr. Mengano de Cual}
% =====================================================================


% ==================== IDIOMA Y CODIFICACIÓN ====================
\usepackage[utf8]{inputenc}
\usepackage[T1]{fontenc}
\usepackage[spanish,es-tabla,es-noshorthands]{babel}
\usepackage{lmodern}
\usepackage{microtype}

% ==================== PÁGINA ====================
\usepackage[letterpaper,margin=2.3cm,headheight=14pt]{geometry}
\usepackage{parskip}
\usepackage{fancyhdr}
\usepackage{enumitem}

% ==================== MATEMÁTICAS ====================
\usepackage{amsmath,amssymb,amsthm,mathtools}

% ==================== FIGURAS, TABLAS, CÓDIGO ====================
\usepackage{graphicx}
\usepackage{booktabs}
\usepackage{array}
% float da el especificador [H], que clava una tabla donde se
% escribio. Una traza que LaTeX manda dos paginas adelante deja
% de leerse junto al paso que la produjo
\usepackage{float}
\usepackage{xcolor}
% tikz ya viene arrastrado por tcolorbox[most], pero se carga
% explicito para no depender de ese detalle; la libreria babel
% protege a tikz de los caracteres activos del español
\usepackage{tikz}
\usetikzlibrary{babel}
% pgfplots dibuja graficas de datos (ejes, escalas log) sobre tikz
\usepackage{pgfplots}
\pgfplotsset{compat=1.18}
\usepackage{listings}
\usepackage[ruled,section]{algorithm}
\usepackage{algpseudocode}
\usepackage[most]{tcolorbox}
% soul da \hl, el unico resaltado que parte renglon; \colorbox no
\usepackage{soul}
\usepackage{hyperref}

% =====================================================================
%  ESTILO
% =====================================================================

% --- encabezado y pie de cada página ---
\pagestyle{fancy}
\fancyhf{}
\fancyhead[L]{\small\textsc{\materia}}
\fancyhead[R]{\small\periodo}
\fancyfoot[C]{\small\thepage}
\renewcommand{\headrulewidth}{0.4pt}

% --- enlaces ---
\hypersetup{
  colorlinks=true,
  linkcolor=black,
  urlcolor=blue,
  citecolor=black,
  pdftitle={Apuntes de \materia},
  pdfauthor={\alumno}
}

% --- paleta ---
\definecolor{acento}{HTML}{1F4E79}
\definecolor{fondoIdea}{HTML}{EAF2FB}
\definecolor{fondoOjo}{HTML}{FDF0E6}
\definecolor{fondoPend}{HTML}{F2F0F7}
\definecolor{comentario}{HTML}{5A7D5A}
\definecolor{cadena}{HTML}{A03030}
\definecolor{fondoCodigo}{HTML}{F7F7F7}
\definecolor{fondoResultado}{HTML}{EAF5EC}
\definecolor{marcador}{HTML}{FFF3A3}

% --- entornos de teorema; numeran por sección, o sea por clase ---
\theoremstyle{definition}
\newtheorem{definicion}{Definición}[section]
\newtheorem{ejemplo}[definicion]{Ejemplo}
\newtheorem{ejercicio}[definicion]{Ejercicio}

\theoremstyle{plain}
\newtheorem{teorema}[definicion]{Teorema}
\newtheorem{lema}[definicion]{Lema}
\newtheorem{corolario}[definicion]{Corolario}
\newtheorem{proposicion}[definicion]{Proposición}

\theoremstyle{remark}
\newtheorem*{observacion}{Observación}
\newtheorem*{quick_sum}{Resumen rápido}

% --- cajas de color para lo que se relee antes del examen ---
\newtcolorbox{idea}[1][]{
  colback=fondoIdea, colframe=acento, boxrule=0.6pt,
  arc=2pt, left=6pt, right=6pt, top=5pt, bottom=5pt,
  title={Idea clave}, fonttitle=\bfseries\small, #1
}
\newtcolorbox{ojo}[1][]{
  colback=fondoOjo, colframe=orange!70!black, boxrule=0.6pt,
  arc=2pt, left=6pt, right=6pt, top=5pt, bottom=5pt,
  title={Ojito...}, fonttitle=\bfseries\small, #1
}
\newtcolorbox{pendiente}[1][]{
  colback=fondoPend, colframe=violet!60!black, boxrule=0.6pt,
  arc=2pt, left=6pt, right=6pt, top=5pt, bottom=5pt,
  title={Pendiente por resolver}, fonttitle=\bfseries\small, #1
}
\newtcolorbox{idear}[1][]{
  colback=fondoPend, colframe=pink!60!black, boxrule=0.6pt,
  arc=2pt, left=6pt, right=6pt, top=5pt, bottom=5pt,
  title={Idea Random}, fonttitle=\bfseries\small, #1
}

% --- el resultado al que llego un ejercicio, para que se vea de
%     lejos al hojear la hoja antes del examen ---
\newtcolorbox{resultado}[1][]{
  colback=fondoResultado, colframe=green!45!black, boxrule=0.6pt,
  arc=2pt, left=6pt, right=6pt, top=5pt, bottom=5pt,
  title={Resultado}, fonttitle=\bfseries\small, #1
}

% --- código ---
\lstdefinestyle{codigo}{
  backgroundcolor=\color{fondoCodigo},
  basicstyle=\ttfamily\footnotesize,
  keywordstyle=\color{acento}\bfseries,
  commentstyle=\color{comentario}\itshape,
  stringstyle=\color{cadena},
  numbers=left, numberstyle=\tiny\color{gray}, numbersep=8pt,
  frame=single, rulecolor=\color{gray!40},
  breaklines=true, showstringspaces=false,
  tabsize=4, captionpos=b,
  extendedchars=true, inputencoding=utf8,
  literate={á}{{\'a}}1 {é}{{\'e}}1 {í}{{\'i}}1
           {ó}{{\'o}}1 {ú}{{\'u}}1 {ñ}{{\~n}}1
           {Ü}{{\"U}}1 {ü}{{\"u}}1 {¿}{{?`}}1
           {¡}{{!`}}1
}
\lstset{style=codigo}
\renewcommand{\lstlistingname}{Código}
\renewcommand{\lstlistlistingname}{Índice de código}
% --- pseudocódigo; algorithmicx arma el cuerpo y algorithm el
%     flotante que lo numera y lo pone en el índice. Las palabras
%     clave vienen en inglés de fábrica y babel no las toca, así que
%     se renombran una por una ---
\floatname{algorithm}{Algoritmo}
\renewcommand{\listalgorithmname}{Índice de algoritmos}

\algrenewcommand\algorithmicrequire{\textbf{Entrada:}}
\algrenewcommand\algorithmicensure{\textbf{Salida:}}
\algrenewcommand\algorithmicend{\textbf{fin}}
\algrenewcommand\algorithmicif{\textbf{si}}
\algrenewcommand\algorithmicthen{\textbf{entonces}}
\algrenewcommand\algorithmicelse{\textbf{si no}}
\algrenewcommand\algorithmicfor{\textbf{para}}
\algrenewcommand\algorithmicforall{\textbf{para todo}}
\algrenewcommand\algorithmicdo{\textbf{hacer}}
\algrenewcommand\algorithmicwhile{\textbf{mientras}}
\algrenewcommand\algorithmicrepeat{\textbf{repetir}}
\algrenewcommand\algorithmicuntil{\textbf{hasta que}}
\algrenewcommand\algorithmicloop{\textbf{ciclo}}
\algrenewcommand\algorithmicprocedure{\textbf{procedimiento}}
\algrenewcommand\algorithmicfunction{\textbf{función}}
\algrenewcommand\algorithmicreturn{\textbf{devolver}}
\algrenewcomment[1]{\hfill{\color{comentario}\itshape // #1}}

% Las dos que algorithmicx no trae. El \ del final es obligatorio: sin
% el TeX se come el espacio y sale "hastan-1" en vez de "hasta n-1"
\algnewcommand\Hasta{\textbf{hasta}\ }
\algnewcommand\Intercambiar{\textbf{intercambiar}\ }

% --- abreviaturas que se usan a cada rato ---
\newcommand{\N}{\mathbb{N}}
\newcommand{\Z}{\mathbb{Z}}
\newcommand{\Q}{\mathbb{Q}}
\newcommand{\R}{\mathbb{R}}
\newcommand{\C}{\mathbb{C}}
\newcommand{\abs}[1]{\left\lvert #1 \right\rvert}
\newcommand{\norm}[1]{\left\lVert #1 \right\rVert}
\newcommand{\set}[1]{\left\{ #1 \right\}}
\newcommand{\Ord}[1]{\mathcal{O}\!\left(#1\right)}
% Mayuscula la notacion grande, minuscula la chica
\newcommand{\Omg}[1]{\Omega\!\left(#1\right)}
\newcommand{\Tht}[1]{\Theta\!\left(#1\right)}
\newcommand{\ord}[1]{o\!\left(#1\right)}
\newcommand{\omg}[1]{\omega\!\left(#1\right)}
% Los tres limites sobre n, que salen en cada renglon del criterio
% del cociente para la notacion asintotica
\newcommand{\LM}{\lim_{n \to \infty}}
\newcommand{\LS}{\limsup_{n \to \infty}}
\newcommand{\LI}{\liminf_{n \to \infty}}

% --- encabezado de cada sesión: título a la izquierda, fecha a la
%     derecha. El argumento corto es el que sale en el índice ---
\newcommand{\clase}[2]{%
  \section[#1]{#1 \hfill \normalfont\small\textit{#2}}%
}

% --- portadilla y tabla de contenidos. Se llama una vez, justo
%     despues de \begin{document}. El renglon del profesor va dentro
%     del center, asi que \profesor puede traer varios nombres
%     separados por \\ y salen apilados y centrados sin tocar nada ---
\newcommand{\portadilla}{%
  \begin{center}
    {\LARGE\bfseries\materia\par}
    \vspace{0.4em}
    {\large \profesor \par}
    \vspace{0.2em}
    {\small \periodo\quad-\quad\alumno\par}
    \vspace{0.3em}
    \rule{\linewidth}{0.5pt}
  \end{center}
  \vspace{0.5em}
  \tableofcontents
  \vspace{1em}
  \rule{\linewidth}{0.5pt}%
}

% =====================================================================
%  DOCUMENTOS DE SOLUCIONES
%  Lo que usa una tanda de ejercicios resueltos. Una libreta no toca
%  nada de aqui, pero vive en el preambulo compartido para no tener
%  dos archivos que arreglar cuando cambie el estilo.
% =====================================================================

% --- encabezado compacto: titulo, subtitulo y regla. Sustituye a
%     \portadilla cuando el documento es una tanda de ejercicios: sin
%     indice, porque se hojea de corrido, no se navega ---
\newcommand{\encabezadoSoluciones}[2]{%
  % El paquete algorithm se carga con la opcion [section] para que en
  % una libreta el pseudocodigo se numere por clase. Aqui no hay
  % secciones, asi que esa numeracion sale como "Algoritmo 0.1". Se
  % corrige al vuelo, y solo en los documentos que llaman a este
  % encabezado: la libreta no se entera.
  \renewcommand{\thealgorithm}{\arabic{algorithm}}%
  \begin{center}
    {\LARGE\bfseries #1\par}
    \vspace{0.3em}
    {\small\itshape #2\par}
    \vspace{0.5em}
    \rule{\linewidth}{0.5pt}
  \end{center}
  \vspace{0.3em}%
}

% --- hoja de un solo ejercicio: un renglon chico arriba con el
%     titulo y su numero, y ya. Sin titulo grande y sin pie de pagina,
%     porque en una hoja de dos paginas cada centimetro que se va en
%     adornos es un paso del desarrollo que no cabe ---
\newcommand{\encabezadoEjercicio}[3]{%
  % Misma correccion que en \encabezadoSoluciones: sin secciones, el
  % pseudocodigo saldria numerado "0.1"
  \renewcommand{\thealgorithm}{\arabic{algorithm}}%
  \fancyhf{}%
  \fancyhead[L]{\small\textsc{#1} (ejercicio #2 de #3)}%
  % El nombre va a la derecha porque las hojas de respuestas se
  % entregan y el profesor pide que cada una lo lleve
  \fancyhead[R]{\small\alumno}%
  \renewcommand{\headrulewidth}{0.4pt}%
  \pagestyle{fancy}%
  \thispagestyle{fancy}%
}

% --- bloque tematico, numerado con romanos: I, II, III. Agrupa los
%     ejercicios que comparten tema. No dibuja regla: la del
%     \problema que sigue ya separa el titulo del primer ejercicio,
%     y dos rayas seguidas dejan un hueco feo ---
\newcounter{parte}
\renewcommand{\theparte}{\Roman{parte}}
\newcommand{\parte}[1]{%
  \bigskip
  \refstepcounter{parte}%
  \noindent{\large\bfseries \theparte. #1\par}%
}

% --- un ejercicio resuelto. Se numera solo y corrido a lo largo del
%     documento, sin reiniciar en cada parte, para poder decir "el 14"
%     sin aclarar de que bloque ---
\newcounter{problema}
\newcommand{\problema}[1]{%
  \bigskip
  \noindent\rule{\linewidth}{0.4pt}\par
  \vspace{0.3em}
  \refstepcounter{problema}%
  \noindent{\bfseries Ejercicio \theproblema\ --- #1\par}
  \vspace{0.3em}%
}

% --- rotulo del paso que sigue: Caso base, Hipotesis inductiva, Paso
%     inductivo. El texto sigue en el mismo renglon ---
\newcommand{\paso}[1]{\par\vspace{0.3em}\noindent\textbf{#1}\ }

% --- la respuesta de un ejercicio, centrada y en un cuadro. El
%     argumento va en modo matematico, sin los delimitadores ---
\newcommand{\respuesta}[1]{%
  \[
    \boxed{\;#1\;}
  \]
}

% --- resaltado inline, como pasarle el plumon a media frase. No
%     admite matematicas adentro: soul rompe con $...$, asi que una
%     formula que hay que destacar va en \respuesta ---
\sethlcolor{marcador}
\newcommand{\resaltar}[1]{\hl{#1}}

% --- justificacion al final de un renglon de \align, en chiquito y
%     entre parentesis, para decir de donde salio el paso ---
\newcommand{\razon}[1]{\quad\text{\small\itshape (#1)}}

%############################ END OF PREAMBULO.TEX #############################
%###############################################################################

%  \begin{document}
%  \portadilla
%
%  Con dos o mas profesores se declara UN solo \profesor y los nombres
%  se separan con \\, uno por renglon:
%
%  \newcommand{\profesor}{Dra. Fulana de Tal \\
%                         Dr. Mengano de Cual}
% =====================================================================


% ==================== IDIOMA Y CODIFICACIÓN ====================
\usepackage[utf8]{inputenc}
\usepackage[T1]{fontenc}
\usepackage[spanish,es-tabla,es-noshorthands]{babel}
\usepackage{lmodern}
\usepackage{microtype}

% ==================== PÁGINA ====================
\usepackage[letterpaper,margin=2.3cm,headheight=14pt]{geometry}
\usepackage{parskip}
\usepackage{fancyhdr}
\usepackage{enumitem}

% ==================== MATEMÁTICAS ====================
\usepackage{amsmath,amssymb,amsthm,mathtools}

% ==================== FIGURAS, TABLAS, CÓDIGO ====================
\usepackage{graphicx}
\usepackage{booktabs}
\usepackage{array}
% float da el especificador [H], que clava una tabla donde se
% escribio. Una traza que LaTeX manda dos paginas adelante deja
% de leerse junto al paso que la produjo
\usepackage{float}
\usepackage{xcolor}
% tikz ya viene arrastrado por tcolorbox[most], pero se carga
% explicito para no depender de ese detalle; la libreria babel
% protege a tikz de los caracteres activos del español
\usepackage{tikz}
\usetikzlibrary{babel}
% pgfplots dibuja graficas de datos (ejes, escalas log) sobre tikz
\usepackage{pgfplots}
\pgfplotsset{compat=1.18}
\usepackage{listings}
\usepackage[ruled,section]{algorithm}
\usepackage{algpseudocode}
\usepackage[most]{tcolorbox}
% soul da \hl, el unico resaltado que parte renglon; \colorbox no
\usepackage{soul}
\usepackage{hyperref}

% =====================================================================
%  ESTILO
% =====================================================================

% --- encabezado y pie de cada página ---
\pagestyle{fancy}
\fancyhf{}
\fancyhead[L]{\small\textsc{\materia}}
\fancyhead[R]{\small\periodo}
\fancyfoot[C]{\small\thepage}
\renewcommand{\headrulewidth}{0.4pt}

% --- enlaces ---
\hypersetup{
  colorlinks=true,
  linkcolor=black,
  urlcolor=blue,
  citecolor=black,
  pdftitle={Apuntes de \materia},
  pdfauthor={\alumno}
}

% --- paleta ---
\definecolor{acento}{HTML}{1F4E79}
\definecolor{fondoIdea}{HTML}{EAF2FB}
\definecolor{fondoOjo}{HTML}{FDF0E6}
\definecolor{fondoPend}{HTML}{F2F0F7}
\definecolor{comentario}{HTML}{5A7D5A}
\definecolor{cadena}{HTML}{A03030}
\definecolor{fondoCodigo}{HTML}{F7F7F7}
\definecolor{fondoResultado}{HTML}{EAF5EC}
\definecolor{marcador}{HTML}{FFF3A3}

% --- entornos de teorema; numeran por sección, o sea por clase ---
\theoremstyle{definition}
\newtheorem{definicion}{Definición}[section]
\newtheorem{ejemplo}[definicion]{Ejemplo}
\newtheorem{ejercicio}[definicion]{Ejercicio}

\theoremstyle{plain}
\newtheorem{teorema}[definicion]{Teorema}
\newtheorem{lema}[definicion]{Lema}
\newtheorem{corolario}[definicion]{Corolario}
\newtheorem{proposicion}[definicion]{Proposición}

\theoremstyle{remark}
\newtheorem*{observacion}{Observación}
\newtheorem*{quick_sum}{Resumen rápido}

% --- cajas de color para lo que se relee antes del examen ---
\newtcolorbox{idea}[1][]{
  colback=fondoIdea, colframe=acento, boxrule=0.6pt,
  arc=2pt, left=6pt, right=6pt, top=5pt, bottom=5pt,
  title={Idea clave}, fonttitle=\bfseries\small, #1
}
\newtcolorbox{ojo}[1][]{
  colback=fondoOjo, colframe=orange!70!black, boxrule=0.6pt,
  arc=2pt, left=6pt, right=6pt, top=5pt, bottom=5pt,
  title={Ojito...}, fonttitle=\bfseries\small, #1
}
\newtcolorbox{pendiente}[1][]{
  colback=fondoPend, colframe=violet!60!black, boxrule=0.6pt,
  arc=2pt, left=6pt, right=6pt, top=5pt, bottom=5pt,
  title={Pendiente por resolver}, fonttitle=\bfseries\small, #1
}
\newtcolorbox{idear}[1][]{
  colback=fondoPend, colframe=pink!60!black, boxrule=0.6pt,
  arc=2pt, left=6pt, right=6pt, top=5pt, bottom=5pt,
  title={Idea Random}, fonttitle=\bfseries\small, #1
}

% --- el resultado al que llego un ejercicio, para que se vea de
%     lejos al hojear la hoja antes del examen ---
\newtcolorbox{resultado}[1][]{
  colback=fondoResultado, colframe=green!45!black, boxrule=0.6pt,
  arc=2pt, left=6pt, right=6pt, top=5pt, bottom=5pt,
  title={Resultado}, fonttitle=\bfseries\small, #1
}

% --- código ---
\lstdefinestyle{codigo}{
  backgroundcolor=\color{fondoCodigo},
  basicstyle=\ttfamily\footnotesize,
  keywordstyle=\color{acento}\bfseries,
  commentstyle=\color{comentario}\itshape,
  stringstyle=\color{cadena},
  numbers=left, numberstyle=\tiny\color{gray}, numbersep=8pt,
  frame=single, rulecolor=\color{gray!40},
  breaklines=true, showstringspaces=false,
  tabsize=4, captionpos=b,
  extendedchars=true, inputencoding=utf8,
  literate={á}{{\'a}}1 {é}{{\'e}}1 {í}{{\'i}}1
           {ó}{{\'o}}1 {ú}{{\'u}}1 {ñ}{{\~n}}1
           {Ü}{{\"U}}1 {ü}{{\"u}}1 {¿}{{?`}}1
           {¡}{{!`}}1
}
\lstset{style=codigo}
\renewcommand{\lstlistingname}{Código}
\renewcommand{\lstlistlistingname}{Índice de código}
% --- pseudocódigo; algorithmicx arma el cuerpo y algorithm el
%     flotante que lo numera y lo pone en el índice. Las palabras
%     clave vienen en inglés de fábrica y babel no las toca, así que
%     se renombran una por una ---
\floatname{algorithm}{Algoritmo}
\renewcommand{\listalgorithmname}{Índice de algoritmos}

\algrenewcommand\algorithmicrequire{\textbf{Entrada:}}
\algrenewcommand\algorithmicensure{\textbf{Salida:}}
\algrenewcommand\algorithmicend{\textbf{fin}}
\algrenewcommand\algorithmicif{\textbf{si}}
\algrenewcommand\algorithmicthen{\textbf{entonces}}
\algrenewcommand\algorithmicelse{\textbf{si no}}
\algrenewcommand\algorithmicfor{\textbf{para}}
\algrenewcommand\algorithmicforall{\textbf{para todo}}
\algrenewcommand\algorithmicdo{\textbf{hacer}}
\algrenewcommand\algorithmicwhile{\textbf{mientras}}
\algrenewcommand\algorithmicrepeat{\textbf{repetir}}
\algrenewcommand\algorithmicuntil{\textbf{hasta que}}
\algrenewcommand\algorithmicloop{\textbf{ciclo}}
\algrenewcommand\algorithmicprocedure{\textbf{procedimiento}}
\algrenewcommand\algorithmicfunction{\textbf{función}}
\algrenewcommand\algorithmicreturn{\textbf{devolver}}
\algrenewcomment[1]{\hfill{\color{comentario}\itshape // #1}}

% Las dos que algorithmicx no trae. El \ del final es obligatorio: sin
% el TeX se come el espacio y sale "hastan-1" en vez de "hasta n-1"
\algnewcommand\Hasta{\textbf{hasta}\ }
\algnewcommand\Intercambiar{\textbf{intercambiar}\ }

% --- abreviaturas que se usan a cada rato ---
\newcommand{\N}{\mathbb{N}}
\newcommand{\Z}{\mathbb{Z}}
\newcommand{\Q}{\mathbb{Q}}
\newcommand{\R}{\mathbb{R}}
\newcommand{\C}{\mathbb{C}}
\newcommand{\abs}[1]{\left\lvert #1 \right\rvert}
\newcommand{\norm}[1]{\left\lVert #1 \right\rVert}
\newcommand{\set}[1]{\left\{ #1 \right\}}
\newcommand{\Ord}[1]{\mathcal{O}\!\left(#1\right)}
% Mayuscula la notacion grande, minuscula la chica
\newcommand{\Omg}[1]{\Omega\!\left(#1\right)}
\newcommand{\Tht}[1]{\Theta\!\left(#1\right)}
\newcommand{\ord}[1]{o\!\left(#1\right)}
\newcommand{\omg}[1]{\omega\!\left(#1\right)}
% Los tres limites sobre n, que salen en cada renglon del criterio
% del cociente para la notacion asintotica
\newcommand{\LM}{\lim_{n \to \infty}}
\newcommand{\LS}{\limsup_{n \to \infty}}
\newcommand{\LI}{\liminf_{n \to \infty}}

% --- encabezado de cada sesión: título a la izquierda, fecha a la
%     derecha. El argumento corto es el que sale en el índice ---
\newcommand{\clase}[2]{%
  \section[#1]{#1 \hfill \normalfont\small\textit{#2}}%
}

% --- portadilla y tabla de contenidos. Se llama una vez, justo
%     despues de \begin{document}. El renglon del profesor va dentro
%     del center, asi que \profesor puede traer varios nombres
%     separados por \\ y salen apilados y centrados sin tocar nada ---
\newcommand{\portadilla}{%
  \begin{center}
    {\LARGE\bfseries\materia\par}
    \vspace{0.4em}
    {\large \profesor \par}
    \vspace{0.2em}
    {\small \periodo\quad-\quad\alumno\par}
    \vspace{0.3em}
    \rule{\linewidth}{0.5pt}
  \end{center}
  \vspace{0.5em}
  \tableofcontents
  \vspace{1em}
  \rule{\linewidth}{0.5pt}%
}

% =====================================================================
%  DOCUMENTOS DE SOLUCIONES
%  Lo que usa una tanda de ejercicios resueltos. Una libreta no toca
%  nada de aqui, pero vive en el preambulo compartido para no tener
%  dos archivos que arreglar cuando cambie el estilo.
% =====================================================================

% --- encabezado compacto: titulo, subtitulo y regla. Sustituye a
%     \portadilla cuando el documento es una tanda de ejercicios: sin
%     indice, porque se hojea de corrido, no se navega ---
\newcommand{\encabezadoSoluciones}[2]{%
  % El paquete algorithm se carga con la opcion [section] para que en
  % una libreta el pseudocodigo se numere por clase. Aqui no hay
  % secciones, asi que esa numeracion sale como "Algoritmo 0.1". Se
  % corrige al vuelo, y solo en los documentos que llaman a este
  % encabezado: la libreta no se entera.
  \renewcommand{\thealgorithm}{\arabic{algorithm}}%
  \begin{center}
    {\LARGE\bfseries #1\par}
    \vspace{0.3em}
    {\small\itshape #2\par}
    \vspace{0.5em}
    \rule{\linewidth}{0.5pt}
  \end{center}
  \vspace{0.3em}%
}

% --- hoja de un solo ejercicio: un renglon chico arriba con el
%     titulo y su numero, y ya. Sin titulo grande y sin pie de pagina,
%     porque en una hoja de dos paginas cada centimetro que se va en
%     adornos es un paso del desarrollo que no cabe ---
\newcommand{\encabezadoEjercicio}[3]{%
  % Misma correccion que en \encabezadoSoluciones: sin secciones, el
  % pseudocodigo saldria numerado "0.1"
  \renewcommand{\thealgorithm}{\arabic{algorithm}}%
  \fancyhf{}%
  \fancyhead[L]{\small\textsc{#1} (ejercicio #2 de #3)}%
  % El nombre va a la derecha porque las hojas de respuestas se
  % entregan y el profesor pide que cada una lo lleve
  \fancyhead[R]{\small\alumno}%
  \renewcommand{\headrulewidth}{0.4pt}%
  \pagestyle{fancy}%
  \thispagestyle{fancy}%
}

% --- bloque tematico, numerado con romanos: I, II, III. Agrupa los
%     ejercicios que comparten tema. No dibuja regla: la del
%     \problema que sigue ya separa el titulo del primer ejercicio,
%     y dos rayas seguidas dejan un hueco feo ---
\newcounter{parte}
\renewcommand{\theparte}{\Roman{parte}}
\newcommand{\parte}[1]{%
  \bigskip
  \refstepcounter{parte}%
  \noindent{\large\bfseries \theparte. #1\par}%
}

% --- un ejercicio resuelto. Se numera solo y corrido a lo largo del
%     documento, sin reiniciar en cada parte, para poder decir "el 14"
%     sin aclarar de que bloque ---
\newcounter{problema}
\newcommand{\problema}[1]{%
  \bigskip
  \noindent\rule{\linewidth}{0.4pt}\par
  \vspace{0.3em}
  \refstepcounter{problema}%
  \noindent{\bfseries Ejercicio \theproblema\ --- #1\par}
  \vspace{0.3em}%
}

% --- rotulo del paso que sigue: Caso base, Hipotesis inductiva, Paso
%     inductivo. El texto sigue en el mismo renglon ---
\newcommand{\paso}[1]{\par\vspace{0.3em}\noindent\textbf{#1}\ }

% --- la respuesta de un ejercicio, centrada y en un cuadro. El
%     argumento va en modo matematico, sin los delimitadores ---
\newcommand{\respuesta}[1]{%
  \[
    \boxed{\;#1\;}
  \]
}

% --- resaltado inline, como pasarle el plumon a media frase. No
%     admite matematicas adentro: soul rompe con $...$, asi que una
%     formula que hay que destacar va en \respuesta ---
\sethlcolor{marcador}
\newcommand{\resaltar}[1]{\hl{#1}}

% --- justificacion al final de un renglon de \align, en chiquito y
%     entre parentesis, para decir de donde salio el paso ---
\newcommand{\razon}[1]{\quad\text{\small\itshape (#1)}}

%############################ END OF PREAMBULO.TEX #############################
%###############################################################################

%  \begin{document}
%  \portadilla
%
%  Con dos o mas profesores se declara UN solo \profesor y los nombres
%  se separan con \\, uno por renglon:
%
%  \newcommand{\profesor}{Dra. Fulana de Tal \\
%                         Dr. Mengano de Cual}
% =====================================================================


% ==================== IDIOMA Y CODIFICACIÓN ====================
\usepackage[utf8]{inputenc}
\usepackage[T1]{fontenc}
\usepackage[spanish,es-tabla,es-noshorthands]{babel}
\usepackage{lmodern}
\usepackage{microtype}

% ==================== PÁGINA ====================
\usepackage[letterpaper,margin=2.3cm,headheight=14pt]{geometry}
\usepackage{parskip}
\usepackage{fancyhdr}
\usepackage{enumitem}

% ==================== MATEMÁTICAS ====================
\usepackage{amsmath,amssymb,amsthm,mathtools}

% ==================== FIGURAS, TABLAS, CÓDIGO ====================
\usepackage{graphicx}
\usepackage{booktabs}
\usepackage{array}
% float da el especificador [H], que clava una tabla donde se
% escribio. Una traza que LaTeX manda dos paginas adelante deja
% de leerse junto al paso que la produjo
\usepackage{float}
\usepackage{xcolor}
% tikz ya viene arrastrado por tcolorbox[most], pero se carga
% explicito para no depender de ese detalle; la libreria babel
% protege a tikz de los caracteres activos del español
\usepackage{tikz}
\usetikzlibrary{babel}
% pgfplots dibuja graficas de datos (ejes, escalas log) sobre tikz
\usepackage{pgfplots}
\pgfplotsset{compat=1.18}
\usepackage{listings}
\usepackage[ruled,section]{algorithm}
\usepackage{algpseudocode}
\usepackage[most]{tcolorbox}
% soul da \hl, el unico resaltado que parte renglon; \colorbox no
\usepackage{soul}
\usepackage{hyperref}

% =====================================================================
%  ESTILO
% =====================================================================

% --- encabezado y pie de cada página ---
\pagestyle{fancy}
\fancyhf{}
\fancyhead[L]{\small\textsc{\materia}}
\fancyhead[R]{\small\periodo}
\fancyfoot[C]{\small\thepage}
\renewcommand{\headrulewidth}{0.4pt}

% --- enlaces ---
\hypersetup{
  colorlinks=true,
  linkcolor=black,
  urlcolor=blue,
  citecolor=black,
  pdftitle={Apuntes de \materia},
  pdfauthor={\alumno}
}

% --- paleta ---
\definecolor{acento}{HTML}{1F4E79}
\definecolor{fondoIdea}{HTML}{EAF2FB}
\definecolor{fondoOjo}{HTML}{FDF0E6}
\definecolor{fondoPend}{HTML}{F2F0F7}
\definecolor{comentario}{HTML}{5A7D5A}
\definecolor{cadena}{HTML}{A03030}
\definecolor{fondoCodigo}{HTML}{F7F7F7}
\definecolor{fondoResultado}{HTML}{EAF5EC}
\definecolor{marcador}{HTML}{FFF3A3}

% --- entornos de teorema; numeran por sección, o sea por clase ---
\theoremstyle{definition}
\newtheorem{definicion}{Definición}[section]
\newtheorem{ejemplo}[definicion]{Ejemplo}
\newtheorem{ejercicio}[definicion]{Ejercicio}

\theoremstyle{plain}
\newtheorem{teorema}[definicion]{Teorema}
\newtheorem{lema}[definicion]{Lema}
\newtheorem{corolario}[definicion]{Corolario}
\newtheorem{proposicion}[definicion]{Proposición}

\theoremstyle{remark}
\newtheorem*{observacion}{Observación}
\newtheorem*{quick_sum}{Resumen rápido}

% --- cajas de color para lo que se relee antes del examen ---
\newtcolorbox{idea}[1][]{
  colback=fondoIdea, colframe=acento, boxrule=0.6pt,
  arc=2pt, left=6pt, right=6pt, top=5pt, bottom=5pt,
  title={Idea clave}, fonttitle=\bfseries\small, #1
}
\newtcolorbox{ojo}[1][]{
  colback=fondoOjo, colframe=orange!70!black, boxrule=0.6pt,
  arc=2pt, left=6pt, right=6pt, top=5pt, bottom=5pt,
  title={Ojito...}, fonttitle=\bfseries\small, #1
}
\newtcolorbox{pendiente}[1][]{
  colback=fondoPend, colframe=violet!60!black, boxrule=0.6pt,
  arc=2pt, left=6pt, right=6pt, top=5pt, bottom=5pt,
  title={Pendiente por resolver}, fonttitle=\bfseries\small, #1
}
\newtcolorbox{idear}[1][]{
  colback=fondoPend, colframe=pink!60!black, boxrule=0.6pt,
  arc=2pt, left=6pt, right=6pt, top=5pt, bottom=5pt,
  title={Idea Random}, fonttitle=\bfseries\small, #1
}

% --- el resultado al que llego un ejercicio, para que se vea de
%     lejos al hojear la hoja antes del examen ---
\newtcolorbox{resultado}[1][]{
  colback=fondoResultado, colframe=green!45!black, boxrule=0.6pt,
  arc=2pt, left=6pt, right=6pt, top=5pt, bottom=5pt,
  title={Resultado}, fonttitle=\bfseries\small, #1
}

% --- código ---
\lstdefinestyle{codigo}{
  backgroundcolor=\color{fondoCodigo},
  basicstyle=\ttfamily\footnotesize,
  keywordstyle=\color{acento}\bfseries,
  commentstyle=\color{comentario}\itshape,
  stringstyle=\color{cadena},
  numbers=left, numberstyle=\tiny\color{gray}, numbersep=8pt,
  frame=single, rulecolor=\color{gray!40},
  breaklines=true, showstringspaces=false,
  tabsize=4, captionpos=b,
  extendedchars=true, inputencoding=utf8,
  literate={á}{{\'a}}1 {é}{{\'e}}1 {í}{{\'i}}1
           {ó}{{\'o}}1 {ú}{{\'u}}1 {ñ}{{\~n}}1
           {Ü}{{\"U}}1 {ü}{{\"u}}1 {¿}{{?`}}1
           {¡}{{!`}}1
}
\lstset{style=codigo}
\renewcommand{\lstlistingname}{Código}
\renewcommand{\lstlistlistingname}{Índice de código}
% --- pseudocódigo; algorithmicx arma el cuerpo y algorithm el
%     flotante que lo numera y lo pone en el índice. Las palabras
%     clave vienen en inglés de fábrica y babel no las toca, así que
%     se renombran una por una ---
\floatname{algorithm}{Algoritmo}
\renewcommand{\listalgorithmname}{Índice de algoritmos}

\algrenewcommand\algorithmicrequire{\textbf{Entrada:}}
\algrenewcommand\algorithmicensure{\textbf{Salida:}}
\algrenewcommand\algorithmicend{\textbf{fin}}
\algrenewcommand\algorithmicif{\textbf{si}}
\algrenewcommand\algorithmicthen{\textbf{entonces}}
\algrenewcommand\algorithmicelse{\textbf{si no}}
\algrenewcommand\algorithmicfor{\textbf{para}}
\algrenewcommand\algorithmicforall{\textbf{para todo}}
\algrenewcommand\algorithmicdo{\textbf{hacer}}
\algrenewcommand\algorithmicwhile{\textbf{mientras}}
\algrenewcommand\algorithmicrepeat{\textbf{repetir}}
\algrenewcommand\algorithmicuntil{\textbf{hasta que}}
\algrenewcommand\algorithmicloop{\textbf{ciclo}}
\algrenewcommand\algorithmicprocedure{\textbf{procedimiento}}
\algrenewcommand\algorithmicfunction{\textbf{función}}
\algrenewcommand\algorithmicreturn{\textbf{devolver}}
\algrenewcomment[1]{\hfill{\color{comentario}\itshape // #1}}

% Las dos que algorithmicx no trae. El \ del final es obligatorio: sin
% el TeX se come el espacio y sale "hastan-1" en vez de "hasta n-1"
\algnewcommand\Hasta{\textbf{hasta}\ }
\algnewcommand\Intercambiar{\textbf{intercambiar}\ }

% --- abreviaturas que se usan a cada rato ---
\newcommand{\N}{\mathbb{N}}
\newcommand{\Z}{\mathbb{Z}}
\newcommand{\Q}{\mathbb{Q}}
\newcommand{\R}{\mathbb{R}}
\newcommand{\C}{\mathbb{C}}
\newcommand{\abs}[1]{\left\lvert #1 \right\rvert}
\newcommand{\norm}[1]{\left\lVert #1 \right\rVert}
\newcommand{\set}[1]{\left\{ #1 \right\}}
\newcommand{\Ord}[1]{\mathcal{O}\!\left(#1\right)}
% Mayuscula la notacion grande, minuscula la chica
\newcommand{\Omg}[1]{\Omega\!\left(#1\right)}
\newcommand{\Tht}[1]{\Theta\!\left(#1\right)}
\newcommand{\ord}[1]{o\!\left(#1\right)}
\newcommand{\omg}[1]{\omega\!\left(#1\right)}
% Los tres limites sobre n, que salen en cada renglon del criterio
% del cociente para la notacion asintotica
\newcommand{\LM}{\lim_{n \to \infty}}
\newcommand{\LS}{\limsup_{n \to \infty}}
\newcommand{\LI}{\liminf_{n \to \infty}}

% --- encabezado de cada sesión: título a la izquierda, fecha a la
%     derecha. El argumento corto es el que sale en el índice ---
\newcommand{\clase}[2]{%
  \section[#1]{#1 \hfill \normalfont\small\textit{#2}}%
}

% --- portadilla y tabla de contenidos. Se llama una vez, justo
%     despues de \begin{document}. El renglon del profesor va dentro
%     del center, asi que \profesor puede traer varios nombres
%     separados por \\ y salen apilados y centrados sin tocar nada ---
\newcommand{\portadilla}{%
  \begin{center}
    {\LARGE\bfseries\materia\par}
    \vspace{0.4em}
    {\large \profesor \par}
    \vspace{0.2em}
    {\small \periodo\quad-\quad\alumno\par}
    \vspace{0.3em}
    \rule{\linewidth}{0.5pt}
  \end{center}
  \vspace{0.5em}
  \tableofcontents
  \vspace{1em}
  \rule{\linewidth}{0.5pt}%
}

% =====================================================================
%  DOCUMENTOS DE SOLUCIONES
%  Lo que usa una tanda de ejercicios resueltos. Una libreta no toca
%  nada de aqui, pero vive en el preambulo compartido para no tener
%  dos archivos que arreglar cuando cambie el estilo.
% =====================================================================

% --- encabezado compacto: titulo, subtitulo y regla. Sustituye a
%     \portadilla cuando el documento es una tanda de ejercicios: sin
%     indice, porque se hojea de corrido, no se navega ---
\newcommand{\encabezadoSoluciones}[2]{%
  % El paquete algorithm se carga con la opcion [section] para que en
  % una libreta el pseudocodigo se numere por clase. Aqui no hay
  % secciones, asi que esa numeracion sale como "Algoritmo 0.1". Se
  % corrige al vuelo, y solo en los documentos que llaman a este
  % encabezado: la libreta no se entera.
  \renewcommand{\thealgorithm}{\arabic{algorithm}}%
  \begin{center}
    {\LARGE\bfseries #1\par}
    \vspace{0.3em}
    {\small\itshape #2\par}
    \vspace{0.5em}
    \rule{\linewidth}{0.5pt}
  \end{center}
  \vspace{0.3em}%
}

% --- hoja de un solo ejercicio: un renglon chico arriba con el
%     titulo y su numero, y ya. Sin titulo grande y sin pie de pagina,
%     porque en una hoja de dos paginas cada centimetro que se va en
%     adornos es un paso del desarrollo que no cabe ---
\newcommand{\encabezadoEjercicio}[3]{%
  % Misma correccion que en \encabezadoSoluciones: sin secciones, el
  % pseudocodigo saldria numerado "0.1"
  \renewcommand{\thealgorithm}{\arabic{algorithm}}%
  \fancyhf{}%
  \fancyhead[L]{\small\textsc{#1} (ejercicio #2 de #3)}%
  % El nombre va a la derecha porque las hojas de respuestas se
  % entregan y el profesor pide que cada una lo lleve
  \fancyhead[R]{\small\alumno}%
  \renewcommand{\headrulewidth}{0.4pt}%
  \pagestyle{fancy}%
  \thispagestyle{fancy}%
}

% --- bloque tematico, numerado con romanos: I, II, III. Agrupa los
%     ejercicios que comparten tema. No dibuja regla: la del
%     \problema que sigue ya separa el titulo del primer ejercicio,
%     y dos rayas seguidas dejan un hueco feo ---
\newcounter{parte}
\renewcommand{\theparte}{\Roman{parte}}
\newcommand{\parte}[1]{%
  \bigskip
  \refstepcounter{parte}%
  \noindent{\large\bfseries \theparte. #1\par}%
}

% --- un ejercicio resuelto. Se numera solo y corrido a lo largo del
%     documento, sin reiniciar en cada parte, para poder decir "el 14"
%     sin aclarar de que bloque ---
\newcounter{problema}
\newcommand{\problema}[1]{%
  \bigskip
  \noindent\rule{\linewidth}{0.4pt}\par
  \vspace{0.3em}
  \refstepcounter{problema}%
  \noindent{\bfseries Ejercicio \theproblema\ --- #1\par}
  \vspace{0.3em}%
}

% --- rotulo del paso que sigue: Caso base, Hipotesis inductiva, Paso
%     inductivo. El texto sigue en el mismo renglon ---
\newcommand{\paso}[1]{\par\vspace{0.3em}\noindent\textbf{#1}\ }

% --- la respuesta de un ejercicio, centrada y en un cuadro. El
%     argumento va en modo matematico, sin los delimitadores ---
\newcommand{\respuesta}[1]{%
  \[
    \boxed{\;#1\;}
  \]
}

% --- resaltado inline, como pasarle el plumon a media frase. No
%     admite matematicas adentro: soul rompe con $...$, asi que una
%     formula que hay que destacar va en \respuesta ---
\sethlcolor{marcador}
\newcommand{\resaltar}[1]{\hl{#1}}

% --- justificacion al final de un renglon de \align, en chiquito y
%     entre parentesis, para decir de donde salio el paso ---
\newcommand{\razon}[1]{\quad\text{\small\itshape (#1)}}

%############################ END OF PREAMBULO.TEX #############################
%###############################################################################


% El preambulo compartido pone el numero de pagina en el pie. Un
% reporte de tres hojas no lo necesita, asi que se limpia solo aqui.
\fancyfoot[C]{}

% =====================================================================
\begin{document}
% =====================================================================

\begin{center}
  {\LARGE\bfseries Actividad 1\par}
  \vspace{0.3em}
  {\large Preprocesamiento y vocabulario de la colección Time\par}
  \vspace{0.5em}
  {\small \alumno \quad-\quad \periodo\par}
  \vspace{0.4em}
  \rule{\linewidth}{0.5pt}
\end{center}

\section*{Resumen}

Se preprocesó la colección Time y se extrajo el vocabulario con su
frecuencia de término, tanto de los 423 documentos como de las 83
consultas, utilizando el texto en crudo, sin palabras nulas,
truncado con Porter original y con la variante de NLTK. 
El preprocesamiento consistió en eliminar las palabras
vacías de la lista propia de la colección y truncar con el algoritmo
de Porter. El vocabulario global pasó de 29\,418 términos en el
texto crudo a 16\,675 tras el pipeline completo, una reducción del
$43.3\%$. También hicieron comparaciones de los vocabularios en las 
diferentes etapas del preprocesamiento.

\section{La colección}

Time es una colección de prueba clásica de recuperación de
información, distribuida por el IR Group de Glasgow. Son artículos de
la revista Time de 1963, en texto ASCII, todo en mayúsculas y con la
puntuación mayormente separada por espacios.

\begin{center}
\begin{tabular}{llr}
    \toprule
    \textbf{Archivo} & \textbf{Contenido} & \textbf{Cantidad} \\
    \midrule
    \texttt{TIME.ALL} & Documentos, marcador \texttt{*TEXT} & 423 \\
    \texttt{TIME.QUE} & Consultas, marcador \texttt{*FIND}  & 83 \\
    \texttt{TIME.REL} & Juicios de relevancia               & 83 \\
    \texttt{TIME.STP} & Palabras vacías                     & 340 \\
    \bottomrule
\end{tabular}
\end{center}

\section{Método}

El procesamiento se implementó en \texttt{actividad.ipynb}, un
\emph{notebook} de Jupyter sobre Python 3.13. Se apoya en dos
bibliotecas y en nada más:

\begin{center}
\begin{tabular}{lll}
    \toprule
    \textbf{Biblioteca} & \textbf{Versión} & \textbf{Para qué} \\
    \midrule
    \texttt{nltk}   & 3.10.3 & El \texttt{PorterStemmer}, en sus dos
                                modos \\
    \texttt{pandas} & 3.0.5  & Las tablas de vocabulario y la matriz
                                documento-término \\
    \bottomrule
\end{tabular}
\end{center}

La lectura, la tokenización y el filtrado de palabras vacías se
resolvieron con la biblioteca estándar, principalmente
\texttt{collections.Counter} para las frecuencias y
\texttt{str.strip} para la puntuación. Las versiones están fijadas en
\texttt{../requirements.txt}, porque el resultado depende de ellas: el
\texttt{PorterStemmer} de NLTK cambia de comportamiento entre modos, y
eso mueve el vocabulario.

El procesamiento son cuatro etapas, aplicadas por igual a documentos y
consultas:

\begin{enumerate}[itemsep=0.15em]
    \item \textbf{Lectura.} Se separa cada archivo por su marcador. Un
          documento va desde su \texttt{*TEXT} hasta el siguiente, o
          hasta \texttt{*STOP} al final.
    \item \textbf{Tokenización.} Se parte por espacios y se quita la
          puntuación de los extremos de cada token, conservando la
          interna para no destrozar siglas como \emph{U.S}.
    \item \textbf{Eliminación de palabras vacías.} Se comparan los
          tokens limpios contra las 340 entradas de
          \texttt{TIME.STP}.
    \item \textbf{Truncamiento.} Porter, en sus dos variantes: el
          algoritmo original de 1980 y el de NLTK por omisión.
\end{enumerate}

El vocabulario se representa en forma larga, una fila por par
documento-término con su frecuencia, y se escribe al formato pedido:
una línea por documento con su identificador seguido de los pares
término y frecuencia.

\section{Resultados}
  
Después de aplicar cada etapa de procesamiento, el vocabulario de la colección
se reduce progresivamente, como se muestra en la Tabla~\ref{tab:docs}.

\begin{table}[H]
  \centering
  \begin{tabular}{lrrrr}
    \toprule
    \textbf{Etapa} & \textbf{Tokens} & \textbf{Vocabulario} &
      \textbf{Pares} & \textbf{Reducción} \\
    \midrule
    Texto crudo        & 263\,749 & 29\,418 & 140\,756 &
      --- \\
    Sin palabras vacías & 124\,030 & 23\,338 & 96\,865 &
      $20.7\%$ \\
    Porter (1980)      & 124\,030 & 16\,675 & 92\,091 &
      $28.6\%$ \\
    Porter (NLTK)      & 124\,030 & 16\,623 & 92\,040 &
      $28.8\%$ \\
    \midrule
    \textbf{Acumulado} & $-53.0\%$ & $-43.3\%$ & $-34.6\%$ & \\
    \bottomrule
  \end{tabular}
  \caption{Efecto de cada etapa sobre los 423 documentos. La reducción
    de cada renglón es sobre el vocabulario del anterior; las dos
    variantes de Porter parten de la misma etapa. El acumulado es con
    Porter de 1980.}
  \label{tab:docs}
\end{table}

Y, al construir la matriz del lenguaje como lo vimos en clase, con los
documentos como instancias y las palabras como atributos, se obtuvo una
matriz dispersa con las frecuencias de cada término en cada documento,
de dimensiones distintas según la etapa del preprocesamiento.

\begin{table}[H]
  \centering
  \begin{tabular}{lrrrr}
    \toprule
    \textbf{Etapa} & \textbf{Dimensiones} & \textbf{Celdas} &
      \textbf{No cero} & \textbf{Dispersión} \\
    \midrule
    Texto crudo         & $423 \times 29\,418$ & 12\,443\,814 &
      140\,756 & $98.87\%$ \\
    Sin palabras vacías & $423 \times 23\,338$ &  9\,871\,974 &
       96\,865 & $99.02\%$ \\
    Porter (1980)       & $423 \times 16\,675$ &  7\,053\,525 &
       92\,091 & $98.69\%$ \\
    Porter (NLTK)       & $423 \times 16\,623$ &  7\,031\,529 &
       92\,040 & $98.69\%$ \\
    \bottomrule
  \end{tabular}
  \caption{Matriz documento-término en cada etapa. Las celdas
    distintas de cero son los pares documento-término que sí
    aparecen.}
  \label{tab:matriz}
\end{table}

Más del $98\%$ de esa matriz son ceros en todas las etapas. En
denso, con enteros de 32 bits, la del texto crudo pesa $49.8$ MB y
la de Porter $28.2$ MB, casi puros ceros. Por eso el vocabulario se
guarda en forma larga, que lista solo lo que sí aparece, y por eso
el formato que pide el enunciado es esa misma forma.

A continuación se muestran los 15 términos más comunes en cada etapa
(crudo, sin vacías, Porter de 1980 y Porter de NLTK), así como las
discrepancias entre los dos modos de Porter.

\begin{table}[H]
  \centering
  \footnotesize
  \begin{tabular}{lr lr lr lr}
    \toprule
    \multicolumn{2}{c}{\textbf{Crudo}} &
    \multicolumn{2}{c}{\textbf{Sin vacías}} &
    \multicolumn{2}{c}{\textbf{Porter 1980}} &
    \multicolumn{2}{c}{\textbf{Porter NLTK}} \\
    \cmidrule(lr){1-2} \cmidrule(lr){3-4}
    \cmidrule(lr){5-6} \cmidrule(lr){7-8}
    \emph{the} & 15\,809 & \emph{u.s} & 726 &
      \emph{u.} & 726 & \emph{u.} & 726 \\
    \emph{.} & 12\,148 & \emph{government} & 583 &
      \emph{govern} & 609 & \emph{govern} & 609 \\
    \emph{of} & 7\,219 & \emph{new} & 549 &
      \emph{new} & 603 & \emph{new} & 549 \\
    \emph{to} & 6\,303 & \emph{party} & 378 &
      \emph{parti} & 434 & \emph{parti} & 434 \\
    \emph{a} & 5\,781 & \emph{de} & 335 &
      \emph{nation} & 426 & \emph{nation} & 426 \\
    \emph{and} & 5\,501 & \emph{minister} & 310 &
      \emph{communist} & 420 & \emph{communist} & 420 \\
    \emph{in} & 5\,244 & \emph{west} & 310 &
      \emph{red} & 377 & \emph{red} & 377 \\
    \emph{\texttt{"}} & 4\,946 & \emph{war} & 280 &
      \emph{minist} & 370 & \emph{minist} & 370 \\
    \emph{that} & 2\,450 & \emph{south} & 280 &
      \emph{de} & 355 & \emph{de} & 355 \\
    \emph{for} & 2\,217 & \emph{communist} & 278 &
      \emph{forc} & 331 & \emph{forc} & 331 \\
    \emph{was} & 2\,134 & \emph{red} & 277 &
      \emph{west} & 310 & \emph{west} & 310 \\
    \emph{with} & 1\,830 & \emph{viet} & 267 &
      \emph{war} & 297 & \emph{war} & 297 \\
    \emph{his} & 1\,810 & \emph{british} & 247 &
      \emph{polit} & 281 & \emph{polit} & 281 \\
    \emph{is} & 1\,777 & \emph{president} & 243 &
      \emph{south} & 280 & \emph{south} & 280 \\
    \emph{he} & 1\,659 & \emph{world} & 240 &
      \emph{countri} & 270 & \emph{countri} & 270 \\
    \bottomrule
  \end{tabular}
  \caption{Los 15 términos más frecuentes en cada etapa, con su
    número de apariciones en toda la colección.}
  \label{tab:top15}
\end{table}


Como se puede observar, en crudo los primeros lugares son artículos, 
preposiciones y el punto,
que con 12\,148 apariciones es el segundo término más frecuente de
la colección, ninguno dice nada del contenido. Pero al quitar las vacías
se puede vislumbrar el tema real, o al menos tratar de inferirlo. Y con Porter
las frecuencias \emph{suben} aunque el vocabulario baje, porque
\emph{govern} absorbe a \emph{government} y \emph{governments}.

\begin{table}[H]
  \centering
  \footnotesize
  \begin{tabular}{lllr}
    \toprule
    \textbf{Palabra} & \textbf{Porter 1980} & \textbf{NLTK} &
      \textbf{Veces} \\
    \midrule
    \emph{days} & \emph{dai} & \emph{day} & 136 \\
    \emph{away} & \emph{awai} & \emph{away} & 100 \\
    \emph{money} & \emph{monei} & \emph{money} & 61 \\
    \emph{trying} & \emph{try} & \emph{tri} & 60 \\
    \emph{news} & \emph{new} & \emph{news} & 54 \\
    \emph{assembly} & \emph{assembli} & \emph{assembl} & 45 \\
    \emph{pay} & \emph{pai} & \emph{pay} & 45 \\
    \emph{stay} & \emph{stai} & \emph{stay} & 42 \\
    \emph{died} & \emph{di} & \emph{die} & 40 \\
    \emph{try} & \emph{try} & \emph{tri} & 37 \\
    \emph{eyes} & \emph{ey} & \emph{eye} & 34 \\
    \emph{possibly} & \emph{possibli} & \emph{possibl} & 31 \\
    \emph{key} & \emph{kei} & \emph{key} & 31 \\
    \emph{spy} & \emph{spy} & \emph{spi} & 26 \\
    \emph{age} & \emph{ag} & \emph{age} & 24 \\
    \bottomrule
  \end{tabular}
  \caption{Las 15 discrepancias más frecuentes entre los dos modos de
    Porter, sobre los 124\,030 tokens procesados.}
  \label{tab:porter}
\end{table}

De la misma manera se procesaron las queries, obteniendo diferentes
representaciones dependiendo de la etapa de procesamiento.

\begin{table}[H]
  \centering
  \begin{tabular}{lrrr}
    \toprule
    \textbf{Etapa} & \textbf{Tokens} & \textbf{Vocabulario} &
      \textbf{Pares} \\
    \midrule
    Texto crudo         & 1\,341 & 552 & 1\,228 \\
    Sin palabras vacías & 722        & 449 & 697 \\
    Porter (1980)       & 722        & 413 & 695 \\
    Porter (NLTK)       & 722        & 413 & 695 \\
    \bottomrule
  \end{tabular}
  \caption{Lo mismo sobre las 83 consultas. Son mucho más cortas:
    entre 3 y 18 palabras tras quitar las vacías, media $8.7$.}
  \label{tab:queries}
\end{table}


La eliminación de palabras vacías es la que se lleva más de la mitad
de los \emph{tokens} ($53\%$), pero apenas un quinto del
vocabulario: quita muchas ocurrencias de pocas palabras distintas.
Porter hace lo contrario, no cambia el número de tokens y encoge el
vocabulario otro $28.6\%$, porque solo agrupa formas.

\section{Conclusiones}

El pipeline reduce el vocabulario a menos de la mitad, de 29\,418
a 16\,675 términos, y las dos etapas contribuyen de forma
distinta, quitar palabras vacías recorta ocurrencias, Porter agrupa
formas.

Quedan dos artefactos sin resolver. Porter se come la \emph{S} de
\emph{U.S} tomándola por plural y la convierte en \texttt{u.}, que
además es el término más frecuente del documento 017, con
15 apariciones. Y el apóstrofo interno sobrevive al truncado de los
extremos, así que \emph{britain} y \emph{britain'} quedan como dos
términos distintos en el mismo documento.

\section*{Entregables}

En \texttt{out/}, una línea por documento o consulta con el formato
\texttt{Doc<id> término frecuencia ...}:

\begin{center}
\begin{tabular}{ll}
    \toprule
    \texttt{vocab\_raw.txt} & Texto crudo \\
    \texttt{vocab\_nn.txt} & Sin palabras vacías \\
    \texttt{vocab\_porter.txt} & Sin vacías y con Porter (1980) \\
    \texttt{vocab\_queries\_*.txt} & Lo mismo para las consultas \\
    \bottomrule
\end{tabular}
\end{center}

El código está en \texttt{actividad.ipynb} y el entorno en
\texttt{../requirements.txt}.

\section*{Referencias}

\begin{enumerate}[label={[\arabic*]}, itemsep=0.3em, leftmargin=*]
    \item \textbf{Colección Time.} IR Group, University of Glasgow.
          Colecciones de prueba para recuperación de información.\\
          \url{http://ir.dcs.gla.ac.uk/resources/test_collections/}\\
          De ahí se descargó \texttt{time.tar.gz}, con
          \texttt{TIME.ALL}, \texttt{TIME.QUE}, \texttt{TIME.REL} y
          \texttt{TIME.STP}.
    \item \textbf{Algoritmo de Porter.} Martin Porter. Página oficial
          del algoritmo, con la especificación y las
          implementaciones.\\
          \url{https://tartarus.org/martin/PorterStemmer/}\\
          Publicado originalmente como \emph{An algorithm for suffix
          stripping}, Program, 14(3), pp. 130--137, 1980.
    \item \textbf{NLTK.} La implementación usada es
          \texttt{nltk.stem.PorterStemmer} 3.10.3, en sus modos
          \texttt{ORIGINAL\_ALGORITHM} y \texttt{NLTK\_EXTENSIONS}.
\end{enumerate}

\appendix

\end{document}

%############################# END OF REPORTE.TEX ##############################
%###############################################################################
